% Generated draft: requires mathematical and source review.
\documentclass[11pt]{article}
\usepackage[margin=25mm]{geometry}
\usepackage{fontspec}
\setmainfont{Latin Modern Roman}
\usepackage{amsmath,amssymb,mathtools,unicode-math}
\setmathfont{Latin Modern Math}
\usepackage{enumitem,needspace,xurl,hyperref}
\setlength{\parindent}{0pt}
\setlength{\parskip}{5pt}
\setlength{\emergencystretch}{3em}
\newcommand{\N}{\mathbb{N}}
\newcommand{\Z}{\mathbb{Z}}
\newcommand{\Q}{\mathbb{Q}}
\newcommand{\R}{\mathbb{R}}
\newcommand{\C}{\mathbb{C}}
\newcommand{\F}{\mathbb{F}}
\newcommand{\A}{\mathbb{A}}
\newcommand{\PP}{\mathbb{P}}
\newcommand{\E}{\mathbb{E}}
\newcommand{\eps}{\varepsilon}
\newcommand{\dd}{\,\mathrm{d}}
\begin{document}
\section*{UnsolvedMath ID: 800385}
\textbf{Source problem code:} catalogue entry
\subsection{Definitions and mathematical statement}
For a field \(K\), write \(\PP^1(K)\) for its projective line of \(K\)-rational points. The supplied record calls \(K\) \emph{Hilbertian} if there do not exist finitely many smooth curves \(C_i/K\) and finite maps \(F_i:C_i\to\PP^1_K\), each of degree at least \(2\), such that
\[
\PP^1(K)=\bigcup_{i=1}^{r}F_i\bigl(C_i(K)\bigr).
\]
The record gives all number fields as examples and states that \(\R\), \(\C\), and \(\Q_p\) are not Hilbertian. The conventions concerning the curves and maps in this characterization are not further specified in the supplied record.

The conjecture attributed to Odoni (1985), reproduced as Conjecture 8 in the supplied 2018 source, is the following: for every Hilbertian field \(K\) and every integer \(d\geq2\), there exist coefficients \(a_0,\ldots,a_{d-1}\in K\) and a point \(P\in K\) such that the monic polynomial
\[
f(X)=X^d+a_{d-1}X^{d-1}+\cdots+a_1X+a_0
\]
satisfies
\[
\iota_K(f,P)=1.
\]
The notation \(\iota_K(f,P)\) is retained exactly from the source record, which does not define it. Consequently, the displayed condition is an unresolved part of the formulation; this draft does not assign it a meaning.

The supplied addendum, attributed to Jones (2015), asserts that when \(K\) is a number field, the polynomial can additionally be chosen in \(\mathcal O_K[X]\), where
\[
\mathcal O_K=\{\alpha\in K:\alpha\text{ satisfies a monic polynomial in }\Z[X]\}.
\]
The requirements \(\deg f=d\), monicity, \(P\in K\), and \(\iota_K(f,P)=1\) are retained in this addendum. Its attribution is preserved without independently assessing its literature status.
\subsection{Short English statement}
For every Hilbertian field \(K\) and integer \(d\geq2\), the source conjectures the existence of a monic degree-\(d\) polynomial \(f\in K[X]\) and \(P\in K\) satisfying \(\iota_K(f,P)=1\). The source does not define \(\iota_K\), so this condition still requires clarification. A supplied addendum allows integral coefficients when \(K\) is a number field.
\subsection{Sources}
\begin{itemize}
\item Local UnsolvedMath registry: permanent ID \(800385\), exact source problem code \texttt{catalogue entry}. The supplied record labels the problem open; this label is source metadata, not an independently verified mathematical assessment.
\item Benedetto, DeMarco, Ingram, Jones et al., \emph{Current Trends and Open Problems in Arithmetic Dynamics} (2018), Conjecture 8. Supplied source URL: \url{https://arxiv.org/abs/1806.04980}. The record reports access on 2026-10-16 and extraction by \texttt{arxiv\_latex}; these are supplied provenance details, not a claim of independent access or verification.
\item Original conjecture attribution supplied by the record: Odoni (1985), reference label \texttt{MR805714}. No further bibliographic details are supplied.
\item Number-field addendum attribution supplied by the record: Jones (2015), reference label \texttt{MR3220023}. No further bibliographic details are supplied.
\end{itemize}
\end{document}
